\chapter{Gravity and the Trade Data}\label{ch:b3:gravity}

The models of the last three chapters say why countries trade. This one asks
how much, and finds one of the most reliable regularities in economics: trade
between two countries rises with the size of both and falls with the distance
between them, like the pull between two masses. The chapter derives that
\emph{\omterm{def:b3:gravity:gravity}{gravity equation}} from the CES demand of \cref{ch:b3:scale-variety-firms},
estimates it on the trade of every pair of countries in the world, and uses it to
measure what borders, languages, agreements and distance cost. It ends with the
long record: trade as a share of output rose, collapsed and rose again.

\begin{openingchart}{Exports from one country to another in 2019 against the
product of their GDPs divided by the distance between them, 28{,}569 country
pairs in 50 groups of equal size (natural logs; trade and GDPs in millions of
dollars, distance in kilometres), with the least-squares line on all pairs.}{CEPII
Gravity database, BACI flows; our computation}
\begin{tikzpicture}
\begin{axis}[omscatter, width=0.62\linewidth, height=0.48\linewidth,
  xlabel={$\ln(\mathit{GDP}_o \times \mathit{GDP}_d / \mathit{distance})$},
  ylabel={$\ln(\text{exports})$}]
  \addplot table[x=lx, y=ly] {figdata/bachelor-3/10-gravity/out/opening_bins.csv};
  \addplot[omThm, thick, no marks, sharp plot] table[x=lx, y=ly] {figdata/bachelor-3/10-gravity/out/opening_fit.csv};
\end{axis}
\end{tikzpicture}
\end{openingchart}

% figdata: opening_stats
Across 28{,}569 pairs, one variable --- the product of the two GDPs over the
distance --- accounts for 62~\% of the variation in log trade, and a 1~\% rise in it
goes with 1.18~\% more trade. Few relations in the social sciences fit so
well with so little.

\begin{recall}
\Cref{ch:b3:panel-data} absorbed unobserved differences with fixed effects, and
\cref{ch:b3:binary-outcomes} estimated by maximum likelihood. CES demand, the
\omterm{def:b3:ricardian-trade:iceberg}{iceberg trade cost} and the trade elasticity come from
\cref{ch:b3:scale-variety-firms} and \cref{ch:b3:ricardian-trade}.
\end{recall}

\begin{notation}[Symbols of this chapter]\label{not:b3:gravity:symbols}
$X_{ij}$ is the value of exports from $i$ to $j$; $Y_i$ is $i$'s output, $E_j$ $j$'s
spending and $Y_w$ the world's; $t_{ij} \ge 1$ is the \omterm{def:b3:ricardian-trade:iceberg}{iceberg trade cost} and $\varepsilon = \sigma
- 1$ the trade elasticity. $\Pi_i$ and $P_j$ are the \omterm{def:b3:gravity:multilateral-resistance}{multilateral resistance} terms
(not profit and not a price level of the whole economy).
\end{notation}

\section{The trade data}\label{sec:b3:gravity:data}

\begin{measure}[What the trade data record]
Customs offices record goods as they cross borders: the exporter's office values
them free on board (FOB), at the border, and the importer's adds the cost of
insurance and freight (CIF). Each flow is therefore declared twice, and the two
declarations rarely agree. The flows used in this chapter (BACI, distributed with
the CEPII Gravity database) reconcile the two into one value, FOB, in thousands
of current dollars. They cover goods, not services; a pair with no flow recorded
may have traded too little to report.
\end{measure}

\omcode{code/bachelor-3/10-gravity/listings/load_trade.py}{1}{16}{One year of the world's bilateral trade.}\label{lst:b3:gravity:load}

\omoutput{code/bachelor-3/10-gravity/outputs/load_trade.txt}{224 countries, and most pairs trade.}\label{out:b3:gravity:load}

The file holds only positive flows. Of the 49{,}952 ordered pairs of the 224
countries, 37.7~\% recorded no trade in 2019, and the five largest flows --- China
and Mexico to the United States, the United States and Canada to each other, China
to Hong Kong --- made 8.9~\% of world trade. % output: load_trade

\section{The gravity equation}\label{sec:b3:gravity:equation}

\begin{definition}[Gravity equation; distance elasticity]\label{def:b3:gravity:gravity}
\index{gravity equation}\index{distance elasticity}
A \emph{gravity equation} relates bilateral trade to the sizes of the two economies
and the costs of trading between them:
$X_{ij} = G\, Y_i^{\beta_1} E_j^{\beta_2} \mathit{dist}_{ij}^{\,\delta} e^{z_{ij}'\gamma}$, where $z_{ij}$ are pair
characteristics (a shared border, language, colonial link, trade agreement). The
\emph{distance elasticity} $\delta$ is the percentage change in trade for a 1~\%
increase in distance.
\end{definition}

\begin{definition}[Multilateral resistance]\label{def:b3:gravity:multilateral-resistance}
\index{multilateral resistance}
The \emph{multilateral resistance} terms $\Pi_i$ and $P_j$ are averages of the
trade costs that exporter $i$ faces to all its markets and importer $j$ faces from
all its suppliers. Trade between two countries depends on their bilateral cost
relative to these averages.
\end{definition}

\begin{proposition}[Gravity from CES demand]\label{prop:b3:gravity:structural}
Let each country make one variety, buyers everywhere have \omterm{def:b3:scale-variety-firms:variety}{CES preferences} with
elasticity $\sigma$, and trade costs be \omterm{def:b3:ricardian-trade:iceberg}{iceberg}. If every country's output is sold
($Y_i = \sum_j X_{ij}$), then
\begin{equation}\label{eq:b3:gr:structural}
  X_{ij} = \frac{Y_i E_j}{Y_w} \Bigl(\frac{t_{ij}}{\Pi_i P_j}\Bigr)^{1-\sigma}, \quad
  \Pi_i^{1-\sigma} = \sum_j \Bigl(\frac{t_{ij}}{P_j}\Bigr)^{1-\sigma}\frac{E_j}{Y_w}, \quad
  P_j^{1-\sigma} = \sum_i \Bigl(\frac{t_{ij}}{\Pi_i}\Bigr)^{1-\sigma}\frac{Y_i}{Y_w} .
\end{equation}
\end{proposition}

\begin{proof}
By the CES demand \eqref{eq:b3:sv:demand}, with $p_i$ the factory price of $i$'s
variety (a taste weight can be folded into it), $j$ spends $X_{ij} = (p_it_{ij}/P_j)^{1-\sigma}E_j$
on it, where $P_j^{1-\sigma} = \sum_i (p_it_{ij})^{1-\sigma}$. Market clearing gives $Y_i = p_i^{1-\sigma}
\sum_j (t_{ij}/P_j)^{1-\sigma}E_j = p_i^{1-\sigma} Y_w\Pi_i^{1-\sigma}$, with $\Pi_i$ as defined. So
$p_i^{1-\sigma} = Y_i/(Y_w\Pi_i^{1-\sigma})$. Substituting into $X_{ij}$ gives the first equation,
and into $P_j^{1-\sigma}$ the third.
\end{proof}

Gravity is thus predicted, with unit elasticities on both GDPs and a \omterm{def:b3:gravity:gravity}{distance
elasticity} equal to $-\varepsilon$ times the elasticity of trade costs to distance. It
also says what simple gravity forgets: two countries far from everyone else ---
Australia and New Zealand --- face high resistance with all partners, so they trade
more with each other than their distance alone predicts.

\begin{omfigure}
\begin{tikzpicture}[econ, x=1cm, y=1cm]
  \foreach \n/\x/\y in {A/0/0, B/1.6/0, C/0.8/1.4} \node[draw, circle, inner sep=2pt, font=\scriptsize] (\n) at (\x,\y) {\n};
  \foreach \n/\x/\y in {D/6/0.2, E/7.4/0.2} \node[draw, circle, inner sep=2pt, font=\scriptsize] (\n) at (\x,\y) {\n};
  \draw[omDef, thick] (A) -- (B); \draw[omDef, thick] (A) -- (C); \draw[omDef, thick] (B) -- (C);
  \draw[omThm, very thick] (D) -- (E);
  \draw[guide] (B) -- (D);
  \node[font=\scriptsize, align=center] at (0.8,-0.8) {close neighbours:\\ low resistance};
  \node[font=\scriptsize, align=center] at (6.7,-0.8) {remote pair:\\ high resistance};
\end{tikzpicture}
\par\smallskip{\small \omterm{def:b3:gravity:multilateral-resistance}{Multilateral resistance}. D and E are as far apart as A and B,
but have no other partner nearby: relative to their alternatives, their bilateral
cost is low, and they trade more with each other than A and B do.}
\end{omfigure}

\section{Estimating gravity on logs}\label{sec:b3:gravity:ols}

Taking logs of the definition gives a linear regression. The listing estimates it
for 2019.

\omcode{code/bachelor-3/10-gravity/listings/gravity_ols.py}{1}{17}{Gravity by OLS on logs.}\label{lst:b3:gravity:ols}

\omoutput{code/bachelor-3/10-gravity/outputs/gravity_ols.txt}{The workhorse, naive version.}\label{out:b3:gravity:ols}

% output: gravity_ols
The fit is good (an $R^2$ of 0.66) and the signs are those of every gravity study:
trade rises with both GDPs, falls with distance (elasticity $-1.21$), and is higher
between neighbours, former colonies and partners in an agreement. But the
regression is the ``naive'' gravity the proposition warns about: it omits the
resistance terms, which are correlated with GDP and distance, and its estimates
are biased.

\section{Multilateral resistance and fixed effects}\label{sec:b3:gravity:fe}

The resistance terms are hard to compute, but in a cross-section they are
exporter and importer constants: $\ln X_{ij} = \ln Y_i - \ln \Pi_i^{1-\sigma} + \ln E_j - \ln
P_j^{1-\sigma} + \dots$. A dummy for each exporter and each importer absorbs them,
together with the GDPs --- the fixed effects of \cref{ch:b3:panel-data}, applied to
the two sides of a flow.

\omcode{code/bachelor-3/10-gravity/listings/gravity_fe.py}{1}{17}{Exporter and importer fixed effects.}\label{lst:b3:gravity:fe}

\omoutput{code/bachelor-3/10-gravity/outputs/gravity_fe.txt}{Structural gravity by OLS.}\label{out:b3:gravity:fe}

% output: gravity_fe
With 440 parameters, the \omterm{def:b3:gravity:gravity}{distance elasticity} becomes $-1.52$: doubling the
distance divides trade by $2^{1.515} = 2.86$. A shared border multiplies trade by
$e^{1.027} = 2.79$, a common official language by $e^{0.825} = 2.28$, and a trade
agreement by $e^{0.464} = 1.59$.

\section{Zeros and Poisson pseudo-maximum likelihood}\label{sec:b3:gravity:ppml}

Two problems remain. Zeros have no logarithm, so OLS on logs drops them. And
the log of a variable is not the variable: if the \omterm{def:b3:gravity:gravity}{gravity equation} holds for the
mean of trade, $E[X_{ij} \mid z] = \exp(z'\beta)$, the error is multiplicative, $X_{ij} =
\exp(z'\beta)\eta_{ij}$ with $E[\eta \mid z] = 1$.

\begin{proposition}[Logs and heteroskedasticity]\label{prop:b3:gravity:jensen}
If $E[\eta \mid z] = 1$ and $\eta$ is not constant, then $E[\ln \eta \mid z] < 0$; if the
distribution of $\eta$ depends on $z$, so in general does $E[\ln \eta \mid z]$, and OLS of $\ln X$
on $z$ is inconsistent for $\beta$.
\end{proposition}

\begin{proof}
The logarithm is strictly concave, so by Jensen's inequality $E[\ln \eta \mid z] < \ln
E[\eta \mid z] = 0$. Then $E[\ln X \mid z] = z'\beta + E[\ln \eta \mid z]$. When the last term
varies with $z$, it is part of the regression function, and OLS estimates $\beta$
plus the slope of that term. For example, if $\ln \eta$ is normal with variance $s^2(z)$
and mean $-s^2(z)/2$, then $E[\ln \eta \mid z] = -s^2(z)/2$: wherever the spread differs, so
does the bias.
\end{proof}

\begin{definition}[Poisson pseudo-maximum-likelihood estimator]\label{def:b3:gravity:ppml}
\index{Poisson pseudo-maximum-likelihood estimator}\index{PPML}
The \emph{Poisson pseudo-maximum-likelihood} (PPML) estimator of $\beta$ in $E[X \mid z] =
\exp(z'\beta)$ solves $\sum_{ij}\bigl(X_{ij} - \exp(z_{ij}'\beta)\bigr)z_{ij} = 0$ --- the first-order
condition of a Poisson likelihood, used without assuming that $X$ is Poisson.
\end{definition}

The estimating equation has mean zero whenever the conditional mean is right,
whatever the variance, so \omterm{def:b3:gravity:ppml}{PPML} is consistent under that condition alone; we
admit the proof. It keeps the zeros, and it weights flows in levels. Standard
errors must be robust, since the Poisson variance is not assumed.

\omcode{code/bachelor-3/10-gravity/listings/gravity_ppml.py}{1}{17}{The same equation by \omterm{def:b3:gravity:ppml}{PPML}.}\label{lst:b3:gravity:ppml}

\omoutput{code/bachelor-3/10-gravity/outputs/gravity_ppml.txt}{Smaller elasticities in levels.}\label{out:b3:gravity:ppml}

% output: gravity_ppml
In levels the \omterm{def:b3:gravity:gravity}{distance elasticity} is $-0.65$, the border coefficient 0.54, the
agreement 0.38, and a common language no longer matters (its coefficient,
$-0.03$, is smaller than its standard error). \omterm{def:b3:gravity:ppml}{PPML} gives most weight to the large
flows, between large and often close economies, and there the costs of
distance and of foreign languages appear smaller. The file has no zeros, so
the difference here comes from the weighting and the heteroskedasticity alone
(\cref{exo:b3:gravity:12}).

\begin{omchart}{The distance elasticity estimated each fourth year, 1996--2020,
with exporter and importer fixed effects: OLS on log positive flows and PPML.}{CEPII
Gravity database, BACI flows; our computation}
\begin{tikzpicture}
\begin{axis}[omtime, width=0.6\linewidth, height=0.4\linewidth, xmin=1994, xmax=2022,
  ymin=-1.8, ymax=-0.3, ylabel={Distance elasticity},
  legend style={at={(0.97,0.5)}, anchor=east},
  legend style={font=\footnotesize}]
  \addplot[omDef, thick, mark=*] table[x=year, y=ols] {figdata/bachelor-3/10-gravity/out/distance_by_year.csv};
  \addplot[omThm, thick, mark=square*] table[x=year, y=ppml] {figdata/bachelor-3/10-gravity/out/distance_by_year.csv};
  \legend{OLS on logs, PPML}
\end{axis}
\end{tikzpicture}
\end{omchart}

% figdata: distance_by_year
Neither estimate shrinks over a quarter-century of cheaper transport and
communication: the OLS elasticity moves from $-1.35$ in 1996 to $-1.54$ in 2012 and
$-1.44$ in 2020, the \omterm{def:b3:gravity:ppml}{PPML} one from $-0.55$ to between $-0.64$ and $-0.67$. This is the
\emph{distance puzzle}. A survey of 1{,}467 estimates from 103 papers found that the
estimated effect of distance rose around the middle of the twentieth century and
has stayed high since. % ledger: b3gr:dh.puzzle
Distance is not only freight: it stands for time, information and trust, which
fall less with technology than shipping costs do.

\begin{omchart}{Share of ordered country pairs with no recorded trade, among the
countries recorded each year, 1996--2020 (percent). Coverage of the early years is
narrower.}{CEPII Gravity database, BACI flows; our computation}
\begin{tikzpicture}
\begin{axis}[omtime, width=0.55\linewidth, height=0.32\linewidth, xmin=1995, xmax=2021,
  ymin=0, ymax=75, ylabel={Percent of pairs}]
  \addplot[omDef, thick, mark=none] table[x=year, y=zero_share] {figdata/bachelor-3/10-gravity/out/zeros.csv};
\end{axis}
\end{tikzpicture}
\end{omchart}

% figdata: zeros
More than a third of all pairs record no trade in every year since 2007 (37.7~\% in
2019). Those zeros carry information: they are the pairs too far, too small or
too costly to trade, the extensive margin of \cref{ch:b3:scale-variety-firms}.
OLS on logs throws them away; \omterm{def:b3:gravity:ppml}{PPML} keeps them when they are in the data.

\section{Borders, languages, agreements}\label{sec:b3:gravity:effects}

Thousands of gravity estimates make a literature large enough for its own
\omterm{def:b3:credible-empirical-work:meta-analysis}{meta-analysis}.

\begin{record}[Gravity's typical coefficients]
Over 2{,}508 estimates from 159 papers, the median \omterm{def:b3:gravity:gravity}{distance elasticity} is $-0.89$
($-1.14$ among studies with fixed effects or equivalent methods); the median
coefficients are 0.97 and 0.85 on the exporter's and the importer's GDP, 0.49 on a
shared border and on a common language, 0.91 on a colonial link, 0.47 on a
regional trade agreement (0.28 with fixed effects) and 0.87 on a common currency.
% ledger: b3gr:hm.t4
\end{record}

The most striking cost is the border itself. A study of 1995 found that, after
controlling for size and distance, Canadian provinces traded 22 times as much
with each other as with United States states. With \omterm{def:b3:gravity:multilateral-resistance}{multilateral resistance}
taken into account, the border cuts trade between the two countries by about
44~\% --- still large, but far below the first estimate, which partly reflected
the small size of Canada's economy. % ledger: b3gr:avw.border

\begin{definition}[Tariff-equivalent trade cost]\label{def:b3:gravity:tariff-equivalent}
\index{tariff-equivalent trade cost}
The \emph{tariff-equivalent} of a trade cost measured by a gravity coefficient $b$
is the ad valorem tariff that would change trade as much. Since trade varies as
$t_{ij}^{-\varepsilon}$, a factor that multiplies trade by $e^{b}$ is equivalent to multiplying
$t_{ij}$ by $e^{-b/\varepsilon}$: a tariff of $e^{-b/\varepsilon} - 1$ when $b < 0$, a tariff cut of $1 -
e^{-b/\varepsilon}$ when $b > 0$.
\end{definition}

The tariff-equivalent needs the trade elasticity $\varepsilon$ from outside the
regression; the median of the estimates identified from tariff changes is 5.03,
the value we use. % ledger: b3gr:hm.t5
Not sharing a border is then worth a tariff of $e^{1.027/5.03} - 1 = 22.7~\%$ (with the
fixed-effects coefficient), and an agreement, by \omterm{def:b3:gravity:ppml}{PPML}, a tariff cut of $1 -
e^{-0.380/5.03} = 7.3~\%$ --- of the order of the tariffs that agreements remove.

\begin{model}{The gravity equation}
\assumes CES demand over countries' goods, \omterm{def:b3:ricardian-trade:iceberg}{iceberg trade costs}, and
market clearing.
\predicts Trade proportional to the product of the GDPs, divided by trade costs
relative to \omterm{def:b3:gravity:multilateral-resistance}{multilateral resistance}; a \omterm{def:b3:gravity:gravity}{distance elasticity} around $-1$.
\tested In thousands of estimates (median \omterm{def:b3:gravity:gravity}{distance elasticity} $-0.89$), and in
this chapter's own fits on 2019 data, where size and distance alone account for 62~\%
of the variation in log trade. % ledger: b3gr:hm.t4
\breaks Zeros; services, which it was not built for; trade costs that change over
time; and a \omterm{def:b3:gravity:gravity}{distance elasticity} that does not fall as transport gets cheaper.
\end{model}

\section{The container and the two waves}\label{sec:b3:gravity:waves}

\begin{omchart}{Trade as a share of GDP, (exports + imports)/2 over GDP, 1870--2020:
unweighted mean over the 13 to 18 advanced economies with data in each
year (percent).}{Jord\`a--Schularick--Taylor Macrohistory Database, release 6; our computation}
\begin{tikzpicture}
\begin{axis}[omtime, width=0.7\linewidth, height=0.38\linewidth, xmin=1870, xmax=2020,
  ymin=0, ymax=42, ylabel={Percent of GDP}]
  \addplot[omDef, thick, mark=none] table[x=year, y=share] {figdata/bachelor-3/10-gravity/out/waves.csv};
\end{axis}
\end{tikzpicture}
\end{omchart}

\begin{record}[The first wave of globalisation and its end]
% figdata: waves_points
Across these economies, trade rose from 17.5~\% of GDP in 1870 to 22.8~\% in 1913,
collapsed in the First World War, and fell to 12.3~\% in 1932, in the Depression. It
was back at 17.9~\% in 1950 and passed its 1913 level only in 1974; it reached 35.0~\%
in 2008 and 37.9~\% in 2019. Globalisation came in two waves, and the first was
undone for half a century.
\end{record}

The second wave rode on falling trade costs: tariffs cut in successive rounds of
negotiation and agreements, and a new way of shipping.

\begin{record}[The container]
Bernhofen, El-Sahli and Kneller used the dates at which countries first adopted
container facilities, and the fact that some goods can be containerised and
others cannot, to estimate the effect of containers on product-level trade from
1962 to 1990. Among developed countries, the cumulative effect over twenty years
was about 700~\%, larger than that of tariff liberalisation (working-paper
estimate). % ledger: b3gr:bek.container
\end{record}

\section{Exercises}\label{sec:b3:gravity:exercises}

\begin{exercise}[$\star$]\label{exo:b3:gravity:1}
In the opening chart, by how much does trade rise, on average, when the product
of GDPs over distance doubles? Where do the binned points depart from the line,
and what might explain it?
\end{exercise}

\begin{exercise}[$\star$]\label{exo:b3:gravity:2}
Pair A has GDPs of 2{,}000 and 500 and is 1{,}000 km apart; pair B has GDPs of
2{,}000 and 1{,}000 and is 4{,}000 km apart. With unit GDP elasticities and a \omterm{def:b3:gravity:gravity}{distance
elasticity} of $-0.9$, which pair trades more, and by what factor?
\end{exercise}

\begin{exercise}[$\star$]\label{exo:b3:gravity:3}
With a \omterm{def:b3:gravity:gravity}{distance elasticity} of $-0.9$, by what percentage does trade change when
distance doubles? When it falls by 10~\%?
\end{exercise}

\begin{exercise}[$\star$]\label{exo:b3:gravity:4}
What observation on one pair of countries would contradict gravity? Why is one
such pair not enough to reject it?
\end{exercise}

\begin{exercise}[$\star\star$]\label{exo:b3:gravity:5}
Add a common official language to the regression of the OLS listing. How do the
coefficients on distance and on a colonial link change, and why?
\end{exercise}

\begin{exercise}[$\star\star$]\label{exo:b3:gravity:6}
Estimate the agreement coefficient with fixed effects by OLS and \omterm{def:b3:gravity:ppml}{PPML} for 2005
and 2019. Convert each into a percentage effect on trade. Is the gap between the
two estimators stable?
\end{exercise}

\begin{exercise}[$\star\star$]\label{exo:b3:gravity:7}
With $\varepsilon = 5$, find the tariff-equivalent of not sharing a border, using the
fixed-effects coefficient 1.027 and the \omterm{def:b3:gravity:ppml}{PPML} coefficient 0.537.
\end{exercise}

\begin{exercise}[$\star\star$]\label{exo:b3:gravity:8}
A new port cuts the ad valorem cost of shipping to one partner from 20~\% to 10~\%.
With $\varepsilon = 5$, by how much does trade with that partner rise, before any change in
\omterm{def:b3:gravity:multilateral-resistance}{multilateral resistance}? Will the general-equilibrium effect be larger or smaller,
and what happens to trade with other partners?
\end{exercise}

\begin{exercise}[$\star\star$]\label{exo:b3:gravity:9}
Why does OLS on logs drop zero flows? If zeros are most common between small and
distant countries, in which direction does dropping them bias the \omterm{def:b3:gravity:gravity}{distance
elasticity}?
\end{exercise}

\begin{exercise}[$\star\star\star$]\label{exo:b3:gravity:10}
Estimate the \omterm{def:b3:gravity:gravity}{distance elasticity} with fixed effects by OLS every fourth year from
1996 to 2020, and test for a linear trend. Is the elasticity falling?
\end{exercise}

\begin{exercise}[$\star\star\star$]\label{exo:b3:gravity:11}
Two countries of equal size $Y$ each make one variety; $t_{ii} = 1$ and $t_{12} = t_{21} =
t$. Use \cref{prop:b3:gravity:structural} to find $X_{12}/X_{11}$ and each country's
share of spending on imports. How does the import share depend on $t$ and $\sigma$?
\end{exercise}

\begin{exercise}[$\star\star\star$]\label{exo:b3:gravity:12}
Simulate $y = e^{1 + x}\eta$ with $x$ uniform on $[0, 1]$, $\ln\eta$ normal with variance $0.5 + x$
and mean $-(0.5 + x)/2$. Compare the slope estimated by OLS on $\ln y$ with the \omterm{def:b3:gravity:ppml}{PPML}
slope over 200 samples of 2{,}000. Explain the OLS result with
\cref{prop:b3:gravity:jensen}.
\end{exercise}

\section{Problem: Brenland Joins a Customs Union}\label{sec:b3:gravity:problem}

\begin{problem}[{Weekend problem --- Brenland joins a customs union}]\label{pb:b3:gravity:union}
An invented world has twenty countries: Brenland, Kestria and eighteen others,
N01 to N18. Brenland and four neighbours form a customs union. Trade follows the
structural gravity of \cref{prop:b3:gravity:structural}, observed with noise that
is larger for small flows, and flows below 0.05 billion dalens go unrecorded.
Of the 380 ordered pairs, 20 are pairs of members and 6 record no trade.
% figdata: union_stats
The union's members and the two largest non-members, with their GDP (billions
of dalens) and distance from Brenland (thousands of km):
% figdata: union_table
\begin{center}\small
\pgfplotstabletypeset[col sep=comma,
  columns/name/.style={column name={Country}, string type, column type=l},
  columns/gdp/.style={column name={GDP}, fixed, precision=0},
  columns/dist/.style={column name={Distance}, fixed, fixed zerofill, precision=2},
  columns/member/.style={column name={Member}, fixed, precision=0},
  every head row/.style={after row=\hline}, every last row/.style={after row=\hline}]
  {figdata/bachelor-3/10-gravity/out/union_table.csv}
\end{center}
Brenland's GDP is 120. Three regressions of exports on log distance and a dummy
for pairs of members, the first also on log GDPs (standard errors in parentheses):
% figdata: union_estimates
\begin{center}\small
\pgfplotstabletypeset[col sep=comma,
  columns={method,pairs,gdp_o,gdp_d,dist,dist_se,union,union_se},
  columns/method/.style={column name={Method}, string type, column type=l},
  columns/pairs/.style={column name={Pairs}, fixed, precision=0},
  columns/gdp_o/.style={column name={GDP$_o$}, fixed, fixed zerofill, precision=2, empty cells with={---}},
  columns/gdp_d/.style={column name={GDP$_d$}, fixed, fixed zerofill, precision=2, empty cells with={---}},
  columns/dist/.style={column name={Distance}, fixed, fixed zerofill, precision=3},
  columns/dist_se/.style={column name={}, fixed, fixed zerofill, precision=3, postproc cell content/.append style={/pgfplots/table/@cell content/.add={(}{)}}},
  columns/union/.style={column name={Union}, fixed, fixed zerofill, precision=3},
  columns/union_se/.style={column name={}, fixed, fixed zerofill, precision=3, postproc cell content/.append style={/pgfplots/table/@cell content/.add={(}{)}}},
  every head row/.style={after row=\hline}, every last row/.style={after row=\hline}]
  {figdata/bachelor-3/10-gravity/out/union_estimates.csv}
\end{center}
With the \omterm{def:b3:gravity:ppml}{PPML} coefficients, the model gives Brenland's trade (exports plus
imports, billions of dalens) without and with the union:
% figdata: union_brenland
\begin{center}\small
\pgfplotstabletypeset[col sep=comma,
  columns/partners/.style={column name={Brenland's trade with}, string type, column type=l},
  columns/before/.style={column name={Without union}, fixed, fixed zerofill, precision=1},
  columns/after/.style={column name={With union}, fixed, fixed zerofill, precision=1},
  columns/change/.style={column name={Change (\%)}, fixed, fixed zerofill, precision=1},
  every head row/.style={after row=\hline}, every last row/.style={after row=\hline}]
  {figdata/bachelor-3/10-gravity/out/union_brenland.csv}
\end{center}

\medskip\noindent\textbf{Part I --- gravity by OLS.}
\begin{enumerate}
\item How many pairs does each regression use, and why does OLS use fewer?
\item By the OLS estimate, by how much does trade fall when distance doubles?
\item Using the OLS coefficients, by what factor should Brenland's exports to
Kestria exceed its exports to N10, setting the union aside?
\item By what percentage does membership raise trade between two members,
according to OLS?
\item The GDP coefficients are well below 1. Which part of the structural equation
does this regression omit, and why does that bias it?
\item Compare the OLS distance coefficient with the median of published
estimates.
\end{enumerate}

\medskip\noindent\textbf{Part II --- multilateral resistance and fixed effects.}
\begin{enumerate}[resume]
\item Why do the exporter and importer dummies make the GDPs unnecessary?
\item How many dummies are needed, with a constant?
\item In this world the true \omterm{def:b3:gravity:gravity}{distance elasticity} is $-1$ and the true union
coefficient 0.5. Is the fixed-effects distance estimate within two standard
errors of the truth?
\item Give a 95~\% confidence interval for the union coefficient by OLS with
fixed effects.
\item N12 is large and very close to Brenland. How does that affect Brenland's
\omterm{def:b3:gravity:multilateral-resistance}{multilateral resistance} and its trade with N10?
\item Why would two remote countries in this world trade more with each other
than their distance predicts?
\end{enumerate}

\medskip\noindent\textbf{Part III --- PPML.}
\begin{enumerate}[resume]
\item Why can \omterm{def:b3:gravity:ppml}{PPML} use all 380 pairs?
\item Give the \omterm{def:b3:gravity:ppml}{PPML} union coefficient's 95~\% interval, and the percentage effect
on trade between two members that it implies.
\item The noise is larger for small flows. Use \cref{prop:b3:gravity:jensen} to
explain why OLS on logs and \omterm{def:b3:gravity:ppml}{PPML} disagree.
\item Which estimator came closest to the true \omterm{def:b3:gravity:gravity}{distance elasticity}, and to the
true union coefficient?
\item With $\varepsilon = 5$, find the tariff cut equivalent to the union by \omterm{def:b3:gravity:ppml}{PPML}.
\item In this world, crossing a border divides the cost term by $e^{1.5}$. Find the
tariff-equivalent of the border.
\end{enumerate}

\medskip\noindent\textbf{Part IV --- trade creation and diversion.}
\begin{enumerate}[resume]
\item Brenland's trade with members rises by less than the \omterm{def:b3:gravity:ppml}{PPML} partial effect.
Why?
\item What happens to Brenland's trade with non-members? Name the effect.
\item Why do Brenland's sales at home fall?
\item By what percentage does Brenland's total trade change?
\item Can gravity alone say whether Brenland gains from the union? What else
would you need?
\item Would the union's effect on Brenland be larger if its members were far
away? Explain.
\item State the result: the union's effect on Brenland's trade with members and
with non-members, in percent.
\end{enumerate}
\end{problem}
